\chapter{Using the World Model in the Offline-to-Online Phase}
\label{chap:online-with-wm}

\section{Motivation}
\label{sec:owm-motivation}

The thesis so far has tested two roles for the LeWM predictor. As a
training-time environment substitute it collapsed in the goal-conditioned
setting for four distinct reasons --- compounding prediction error,
$Q$-function corruption by WM-only transitions, threshold pathologies
induced by predictor fine-tuning, and the model-exploitation trap when
the predictor was trained against a value loss
(Chapter~\ref{chap:wm-as-env}). As a subgoal grounder it worked when the
WM was only used at evaluation time --- HIQL+LeWM lifted the offline
baseline modestly (\S\ref{sec:hiql-lewm}) --- but every attempt to feed
WM rollouts back into the training loop, including WGSP
(\S\ref{sec:wgsp}), reproduced the same instability.

We turn next to the setting that originally motivated this thesis:
\emph{offline-to-online} reinforcement learning, in which a policy
pretrained on a fixed dataset is improved during a second
fine-tuning phase
(\S\ref{sec:offline-online-bg}). The canonical move is to collect fresh
transitions in the real environment during that second phase; here we
ask whether the LeWM can stand in for that real interaction. The chapter
opens with a faithful test of that question, framed as a sequence of
WM-as-environment experiments on the offline-to-online benchmark
\texttt{visual-cube-single-singletask-task1-v0}, on which the OGBench
IMPALA baseline reaches $92.8\%$ \cite{espeholt2018impala, ogbench_park2025}
and our offline JEPA policy plateaus at $81.6\%$. We move from the
goal-conditioned variant to the single-task variant because the
offline-to-online literature uses single-task benchmarks against this
IMPALA reference; \S\ref{sec:owm-gc-validation} returns briefly to the
goal-conditioned variant to verify the final method does not collapse
there.

\paragraph{The hypothesis.}
The Chapter~\ref{chap:wm-as-env} synthesis suggests one surviving claim:
if the predictor is accurate enough for \emph{short-horizon} planning ---
which the drift curve of \S\ref{sec:wm-drift} indicates it is, up to
about $\Delta\approx 10$ raw frames --- then it may be usable in the
offline-to-online phase either (a) as the training-time environment, or
(b) as an inference-time planner over $H\le 2$ chunked steps. This
chapter tests both. Option (a) fails in §\ref{sec:owm-exp1}--§\ref{sec:owm-exp8},
recapitulating the same four failure modes in the new setting; option
(b) succeeds in §\ref{sec:mpc-bon}--§\ref{sec:mpc-interleaved} and is
the source of the headline result.

\paragraph{The precise sense in which we do and do not use the WM at
training time.} The final method --- interleaved MPC distillation with
FMQ labels (§\ref{sec:mpc-interleaved}) --- updates the actor on
$(z, \mathbf{a}^{\text{MPC}})$ pairs where $z$ is a real latent and only
the action label is generated by MPC. The critic is never trained on a
WM-imagined next state: its TD target uses real $(z, \mathbf{a}, r, z')$
tuples from the offline buffer. The WM does appear in the actor's
training signal indirectly, through the MPC procedure that selects
$\mathbf{a}^{\text{MPC}}$. This is the precise structural distinction
that separates the working method from every WM-as-environment attempt
in this chapter and Chapter~\ref{chap:wm-as-env}: the critic's
calibration is anchored to real transitions throughout.

\paragraph{Relation to prior work.}
Several established methods sit on either side of the same line ---
whether the WM appears in the critic's training signal.
\begin{itemize}
    \item The LeWM architecture itself ships with a CEM-style planner used
    at evaluation only \cite{maes2026leworldmodelstableendtoendjointembedding}.
    Our best-of-$N$ procedure (§\ref{sec:mpc-bon}) is a learned-proposal
    variant of that planner.
    \item TD-MPC \cite{hansen2022tdmpc} and TD-MPC2 \cite{hansen2024tdmpc2}
    use a learned model and a learned value function jointly in a
    sampling-MPC controller --- the same structural idea as our
    BoN/Grad-MPC family, with sampling-and-elite-selection replaced by
    gradient refinement.
    \item MOPO \cite{yu2020mopomodelbasedofflinepolicy} and
    COMBO \cite{yu2021combo} use the WM directly in the TD target with a
    pessimism penalty. §\ref{sec:owm-exp45} reports that on this task
    even the strongest version of that idea collapses to $2.4\%$, so we
    deliberately do not put the WM in the TD target.
    \item Dreamer \cite{hafner2020dreamer} and SVG-style methods
    \cite{heess2015svg} use the WM as a differentiable simulator and
    backpropagate analytic policy gradients through it.
    §\ref{sec:owm-exp8} shows this collapses on our task.
\end{itemize}

% ---------------------------------------------------------------------
\section{The latents are sufficient: a positive control}
\label{sec:owm-latent-obs}

Before we attribute any failure in this chapter to the world model, we
have to rule out the alternative explanation that the LeJEPA encoder
$E_\xi$ simply does not retain enough of the scene for an offline-to-online
agent to solve the task. The cleanest test is to take everything in our
pipeline except the predictor and let the agent learn against true
dynamics in the real environment. We refer to this control as
\emph{Latent-Observation Online RL}: the encoder is in the loop, the
predictor is not.

\paragraph{Setup.}
A thin observation wrapper sits between the OGBench env and the agent.
Each pixel observation from the real env is passed through the frozen
LeJEPA encoder $E_\xi$ and only its $192$-dimensional latent
$z_t = E_\xi(o_t)$ is exposed to the agent. The dynamics, rewards, and
termination conditions are entirely real. The agent itself is the
action-chunked Flow Q-Learning agent of \S\ref{sec:fql-bg}, with $h=5$
chunked primitive actions. Training runs $5\!\times\!10^5$ offline
steps on the JEPA-encoded play dataset followed by $5\!\times\!10^5$
steps of real-env online interaction.

\paragraph{Result.}
The agent reaches a success rate of $1.00$ at the end of the online
phase ($50$ episodes, random initial states).

\paragraph{What this rules out.}
The $192$-D LeJEPA latents carry enough task-relevant information for an
offline-to-online actor-critic to solve the task perfectly when paired
with true dynamics. Every failure in the rest of the chapter is
therefore attributable to the predictor $f_\psi$, not the encoder
$E_\xi$.

% ---------------------------------------------------------------------
\section{Common setup for the WM-environment experiments}
\label{sec:owm-setup}

The agent and offline data are shared across §\ref{sec:owm-exp1}--§\ref{sec:owm-exp8}.
The agent is the action-chunked FQL agent introduced in \S\ref{sec:fql-bg};
the action chunks are $25$-dimensional (five primitive actions of five
dimensions each), $\gamma = 0.99$, target Polyak rate $\tau = 0.005$,
ensemble of two critics aggregated by mean, BC coefficient $\alpha=100$
on the one-step distillation head. The chunk length $h=5$ is fixed by
the world model, which advances the latent by five primitive actions per
WM step. What varies between experiments is the source of the
transitions $(z, \mathbf{a}, r, z')$ during the online phase.

The WM environment $\widetilde{\mathcal{E}}$ wraps the frozen LeWM
predictor as a Gymnasium env with $192$-dim latent state, $25$-dim
chunked action, $\tilde T = 40$ WM steps per episode ($\approx 200$
real-env steps), and one of two reward shapes:
\begin{align}
    \tilde r_t^{\text{sparse}} &= \mathbb{1}\!\left[\,\|z_t - z_g\|_2 < \tau_{\text{done}}\,\right], \\
    \tilde r_t^{\text{dense}}  &= -\|z_t - z_g\|_2 / s, \qquad s = 10.
\end{align}
Whenever the online phase uses $\tilde r^{\text{dense}}$, the offline
replay buffer is relabelled with the same dense function so the critic
does not have to fit two reward distributions simultaneously.

% ---------------------------------------------------------------------
\section{Plain WM-environment training}
\label{sec:owm-exp1}

The simplest possible setup: replace the real OGBench env entirely with
$\widetilde{\mathcal{E}}$, use sparse rewards, and let the agent run.
Each WM episode starts from a uniformly sampled latent in an
initial-state pool drawn from the offline data and runs for at most
$40$ WM steps with $\tau_{\text{done}}=14$. The agent (starting from a
$61.6\%$ offline checkpoint trained with the unfinetuned LeJEPA WM)
collects synthetic transitions and updates with the standard FQL
losses.

\paragraph{Hypothesis.}
The WM environment supplies an essentially unbounded set of fresh
on-policy transitions in latent space. If the predictor is faithful
enough, the additional coverage should let the policy continue improving
in the same way that online interaction with the real env would.

\paragraph{Result.}
Real-env success rate after $5\!\times\!10^5$ online steps:
$20$--$30\%$ across three seeds.

\paragraph{Failure analysis.}
The WM was trained with a one-step prediction loss, and its per-step
error is small ($\sim\!8$--$9\%$ cosine distance in the play data) but
unbiased --- nothing in the loss keeps a rollout near the offline
support. Over $40$ WM steps the cumulative drift takes the agent's
latent into regions the predictor was never trained to model, and the
critic, fit to TD targets in those drifted regions, learns to
extrapolate Q-values into hallucinated states. The actor then exploits
the extrapolation by selecting actions that maximise the corrupted
critic. The same agent paired with true dynamics
(§\ref{sec:owm-latent-obs}) reaches $100\%$; the difference is entirely
in the predictor. This is the same compounding-error failure mode as in
\S\ref{sec:wm-as-env-synthesis}, recapitulated in the
offline-to-online setting.

% ---------------------------------------------------------------------
\section{Anchored restart rollouts}
\label{sec:owm-exp2}

Compounding error is the most direct symptom of §\ref{sec:owm-exp1};
the most direct fix is to cap the length of any unconstrained rollout.
We do this with an \emph{anchored restart rollout}: every $L_b$ WM
steps the current latent is matched against the offline pool, and if
the match is close enough the rollout is truncated and restarted from
the anchor.

\paragraph{Anchor mechanism.}
We maintain an offline latent pool $\mathcal{Z}_{\text{off}}$ and query
it with FAISS (\S\ref{sec:faiss-bg}) for the nearest neighbour of the
current latent. Define
\begin{equation}
    \Pi^{\epsilon}_{\mathcal{Z}_{\text{off}}}(z) \;=\;
    \begin{cases}
        \hat z & \text{if } \cos(z,\hat z)\ge \tau_{\cos}
        \text{ and } \|z-\hat z\|_2 \le \tau_{L_2},\\
        z & \text{otherwise,}
    \end{cases}
    \label{eq:owm-anchor-proj}
\end{equation}
applied every $L_b = 3$ WM steps with $\tau_{\cos} = 0.95$ and
$\tau_{L_2} = 4.0$. If the nearest neighbour passes the threshold the
rollout is truncated and the next \texttt{reset()} starts from $\hat z$;
otherwise the rollout is abandoned and a fresh latent is drawn. The
transition tuple appended to the replay buffer always uses the raw WM
output as the next state --- there is no fake snap edge in the learned
trajectory.

\paragraph{Three variants and their results.}
\begin{itemize}
    \item Sparse-reward, base WM, $\tau_{\text{done}}=10$, from the
    $61.6\%$ checkpoint: $13$--$22\%$ across four seeds.
    \item Dense-reward, partially fine-tuned WM, $\tau_{\text{done}}=15$:
    $8$--$54\%$ across four seeds (very high variance).
    \item Dense-reward, fully fine-tuned WM, $\tau_{\text{done}}=15$,
    from the $81.6\%$ checkpoint: $15.2\%$.
\end{itemize}

\paragraph{Failure analysis.}
Anchoring bounds the per-branch rollout but the critic is still trained
on sparse-reward signals computed against the WM's notion of ``goal
reached''. That notion does not reliably agree with the real env's: a
rollout that the WM thinks has reached the goal --- and therefore
receives a reward of $1$ --- may correspond to a real-env trajectory
that has not. The full-fine-tune variant fails for a more specific
reason. After fine-tuning, the typical agent-to-goal latent distance
drops from $\approx 25$ with the base WM to $\approx 11$ with the
fully fine-tuned WM. With $\tau_{\text{done}}=15$ --- a threshold
calibrated for the base WM --- every starting latent is already inside
the success region, so every WM episode terminates at step $1$ with the
success flag set. The agent is supervised throughout the entire
$5\!\times\!10^5$ online steps on degenerate one-step episodes carrying
no horizon, no credit assignment, and no exploration; the offline
policy is actively corrupted by training on this signal, and the
$81.6\%$ baseline drops to $15.2\%$. This is the
\emph{termination-threshold pathology}: any reward shaping built on the
WM's latent geometry is sensitive to its calibration, and our
fine-tuning pipeline produces predictors whose geometries differ
enough to make uniform thresholds unsafe.

% ---------------------------------------------------------------------
\section{Joint training of the world model and policy}
\label{sec:owm-exp3}

A more ambitious idea: rather than freeze the WM, update it jointly
with the policy during the offline phase so that the predictor is
encouraged to produce predictions that are useful for $Q$-learning
rather than merely physically plausible.

\paragraph{The four world models in this chapter.}
We use four predictors, all built from the same LeJEPA architecture
but trained on different data and with different gradient pathways:
\begin{itemize}
    \item \emph{Base WM}: trained on the OGBench
    \texttt{visual-cube-single-play-v0} dataset (undirected play; never
    seen the offline task-relabelled dataset).
    \item \emph{Predictor-finetuned WM}: the base WM with only its
    predictor head fine-tuned on the offline play dataset (encoder
    frozen).
    \item \emph{Fully fine-tuned WM}: the base WM with both encoder and
    predictor fine-tuned on the offline play dataset. Used by all the
    fine-tuned experiments in this chapter.
    \item \emph{Joint-trained WM}: the WM in this section, trained
    alongside the policy and critic during the offline phase with the
    value-aware loss described below.
\end{itemize}

\paragraph{Joint-training loss.}
The joint-trained WM has its own AdamW optimiser and receives gradient
from a two-term loss applied during the offline phase:
\begin{align}
    \mathcal{L}_{\text{WM}}(\psi) &= \alpha\,\mathcal{L}_{\text{pred}}(\psi) + \beta(t)\,\mathcal{L}_{\text{value}}(\psi),
    \label{eq:owm-joint}\\
    \mathcal{L}_{\text{pred}}(\psi) &= \mathbb{E}\!\left[\,
        \tfrac{\|f_\psi(z, a) - z'\|_2^2}{\|z'\|_2^2 + \varepsilon}\,\right],
    \label{eq:owm-pred}\\
    \mathcal{L}_{\text{value}}(\psi) &= \mathbb{E}\!\left[\Big(
        Q_{\bar\theta}(z, a) - \big(r + \gamma\, m\, Q_{\bar\theta}(f_\psi(z, a),\, a')\big)\Big)^2\right],
    \quad a' \sim \pi_\theta\big(\,\cdot\mid \mathrm{sg}[f_\psi(z, a)]\big).
    \label{eq:owm-value}
\end{align}
The intent of $\mathcal{L}_{\text{value}}$ is to encourage the WM to
produce predictions that are \emph{self-consistent} with the critic
--- predicted next states whose target-critic value satisfies the
Bellman equation against the real reward. The stop-gradient
$\mathrm{sg}[\cdot]$ keeps the policy from receiving updates through
the WM; the target critic $Q_{\bar\theta}$ keeps the critic itself from
becoming a moving target. The relative weight $\beta(t)$ ramps linearly
from $0$ to $\beta_{\max} = 0.1$ over the first $50\,000$ steps;
$\alpha = 1$.

\paragraph{Result.}
The joint policy reaches $79.6\%$ offline; an anchored dense-reward
online phase that uses the joint WM as the environment and the joint
policy as the starting point yields $9.6\%$.

\paragraph{Diagnostic decomposition.}
Two follow-up experiments isolate the failure by pairing each joint
component with its non-joint counterpart:
\begin{center}
\begin{tabular}{lll r}
\toprule
Run & Policy & WM & SR (\%) \\
\midrule
base policy + joint WM  & base ($81.6\%$) & joint WM &  $7.2$ \\
joint policy + base WM  & joint ($79.6\%$) & base WM & $30.0$ \\
\bottomrule
\end{tabular}
\end{center}
The joint policy paired with the base WM is no worse than other failed
runs in this chapter. The joint WM destroys performance regardless of
which policy is paired with it. The locus of the failure is the
predictor.

\paragraph{Failure analysis: model exploitation.}
$\mathcal{L}_{\text{value}}$ in Eq.~\eqref{eq:owm-value} measures the
Bellman residual between the target critic at the real next state and
the target critic at the WM-predicted next state. Minimising it gives
the predictor a strong incentive to produce $z_{\text{pred}}$ values
that satisfy the Bellman equation under the current critic, regardless
of whether they correspond to plausible real dynamics. The only
counterforce is $\mathcal{L}_{\text{pred}}$. With $\alpha=1$ and
$\beta=0.1$ at the end of the ramp, the value term is a $10\times$
smaller pull, but it acts in a direction that is much easier to
optimise (predicting a single high-value latent vs.\ predicting the
correct latent for arbitrary inputs). MOPO
\cite{yu2020mopomodelbasedofflinepolicy} and COMBO
\cite{yu2021combo} address exactly this concern by freezing the model
and penalising the TD target where the model is uncertain; the next
section tests that family of methods on our problem.

% ---------------------------------------------------------------------
\section{Uncertainty-penalised rewards}
\label{sec:owm-exp45}

A standard remedy for model exploitation is to keep the WM frozen but
penalise predicted transitions by an estimate of model uncertainty
\cite{yu2020mopomodelbasedofflinepolicy, yu2021combo}. We test two
variants inside the anchored, dense-reward, fully fine-tuned-WM
pipeline, both starting from the $81.6\%$ baseline. The penalised
reward is
\begin{equation}
    \tilde r'_t \;=\; \tilde r_t \;-\; \lambda\, u(z_{t+1}, a_t).
    \label{eq:owm-mopo-reward}
\end{equation}

\paragraph{Ensemble disagreement.}
Three independently seeded copies of the fully fine-tuned WM are run in
parallel at every step; per-step uncertainty is the $\ell_2$ norm of
their per-element standard deviation:
\begin{equation}
    u_{\text{ens}}(z, a) \;=\; \Big\|\,\sigma_k\!\big[f_\psi^{(k)}(z, a)\big]\,\Big\|_2,
    \label{eq:owm-mopo-ens}
\end{equation}
with $\lambda = 0.6$. Result: $21.2\%$.

\paragraph{Nearest-neighbour cosine distance.}
A single WM, with uncertainty measured as cosine distance from each
predicted latent to its nearest neighbour in the offline pool:
\begin{equation}
    u_{\text{nn}}(z') \;=\; 1 \;-\; \max_{z\in\mathcal{Z}_{\text{off}}}\;
        \tfrac{\langle z, z'\rangle}{\|z\|\,\|z'\|},
    \label{eq:owm-mopo-nn}
\end{equation}
with $\lambda = 0.07$. Result: $2.4\%$.

\paragraph{Failure analysis.}
The NN-distance result is striking because it is worse than the
unpenalised baseline, and the reason is structural rather than a tuning
issue. The play data was collected by undirected exploration of the
workspace, so the latents it produces are concentrated in the
workspace's interior rather than along the goal-directed trajectories
that solve the task. An NN-distance penalty measures \emph{how far
this state is from the play prior}, not how unreliable the WM's
prediction at this state actually is. Successful task-solving actions
push the cube toward the goal latent --- which is exactly the part of
the latent space the play data does not cover --- so the penalty is
largest precisely on the actions we want the policy to take. The
ensemble variant escapes this specific trap because all three
predictors share the play-data bias, so their disagreement does not
point in the same direction as the bias; but neither penalty
discriminates between WM error and policy correctness in a useful way,
and both runs additionally inherit the termination-threshold pathology
of §\ref{sec:owm-exp2}.

% ---------------------------------------------------------------------
\section{Differentiable WM rollouts}
\label{sec:owm-exp8}

A final option is to abandon the WM-as-environment framing and use the
predictor as a differentiable simulator, backpropagating an analytic
policy gradient through a short imagined rollout into the actor
parameters \cite{heess2015svg, hafner2020dreamer}.

\paragraph{Rollout loss.}
The actor receives an additional term
\begin{equation}
    \mathcal{L}_{\text{roll}}(\theta) \;=\;
        -\mathbb{E}_{z\sim\mathcal{D}}\!\left[\,\sum_{k=0}^{H-1}\gamma^{k}\,
            \tilde r^{\text{dense}}\!\big(z_{t+k+1}\big)\right],
    \quad
    z_{t+k+1} = f_\psi\!\big(z_{t+k},\,\pi_\theta(z_{t+k})\big),
    \label{eq:owm-rollout-loss}
\end{equation}
with $H = 3$. Both $f_\psi$ and the flow-matching actor are
differentiable; the gradient of the dense reward at the rollout
endpoint flows backward through three WM steps and through the actor's
$S = 10$ Euler integration steps into the actor parameters.

\paragraph{Result and failure analysis.}
For both $\lambda_{\text{roll}} \in \{0.1, 1.0\}$, real-env success
rate collapses to $0.0\%$. The gradient passes through $S + H = 13$
noisy compositions before reaching the dense reward, each contributing
variance. The chain is too long for the analytic gradient to point
reliably in the direction of higher real-env return. At
$\lambda_{\text{roll}} = 1$ the rollout term dominates the BC loss and
the actor abandons the behavioural distribution; at $0.1$ the rollout
gradient is too small to influence the actor against the BC term.

% ---------------------------------------------------------------------
\section{Synthesis of the WM-as-environment failures}
\label{sec:owm-synthesis}

Table~\ref{tab:owm-summary} collects every WM-as-environment attempt
of this chapter. None beats its corresponding offline baseline.

\begin{table}[h]
\centering
\caption{Summary of all WM-based offline-to-online attempts on
   \texttt{visual-cube-single-singletask-task1-v0}. SR is the real-env
   success rate over $250$ episodes. The latent-obs control row
   (§\ref{sec:owm-latent-obs}) is included for reference: with the real
   env in place of the WM, the same agent reaches $100\%$.}
\label{tab:owm-summary}
\begin{tabular}{l l c c}
\toprule
Experiment & Mechanism & Offline SR (\%) & Online SR (\%) \\
\midrule
Latent-Obs Online RL (§\ref{sec:owm-latent-obs}) & encoder-only control & 81.6 & \textbf{100.0} \\
\midrule
Plain WM-env, sparse (§\ref{sec:owm-exp1})       & no anchoring          & 61.6 & 20--30 \\
Anchored, sparse (§\ref{sec:owm-exp2})           & $L_b{=}3$, $\tau_{\cos}{=}0.95$ & 61.6 & 13--22 \\
Anchored, dense, partial FT (§\ref{sec:owm-exp2})& dense reward          & 61.6 & 8--54 \\
Anchored, dense, full FT (§\ref{sec:owm-exp2})   & $1$-step degenerate episodes & 81.6 & 15.2 \\
Joint training, anchored (§\ref{sec:owm-exp3})   & $\mathcal{L}_{\text{pred}}+\beta\mathcal{L}_{\text{value}}$ & 79.6 & 9.6 \\
\quad diag: base policy + joint WM               & isolate WM            & 81.6 & 7.2 \\
\quad diag: joint policy + base WM               & isolate policy        & 79.6 & 30.0 \\
Ensemble penalty (§\ref{sec:owm-exp45})          & $\lambda{=}0.6$, $3$ WMs & 81.6 & 21.2 \\
NN-distance penalty (§\ref{sec:owm-exp45})       & $\lambda{=}0.07$       & 81.6 & 2.4 \\
Differentiable rollouts (§\ref{sec:owm-exp8})    & frozen WM, actor BPTT  & 61.6 & 0.0 \\
\bottomrule
\end{tabular}
\end{table}

The encoder-only control rules out the encoder, so every failure above
is attributable to the world-model predictor. The diagnostic chain
identifies four distinct routes by which it goes wrong --- compounding
latent-prediction error (§\ref{sec:owm-exp1}, §\ref{sec:owm-exp2});
$Q$-function corruption (every online run, regardless of anchoring or
pessimism); the termination-threshold pathology
(§\ref{sec:owm-exp2}, §\ref{sec:owm-exp45}); and the
model-exploitation trap when the WM is updated against a
critic-derived loss (§\ref{sec:owm-exp3}). These are the same four
mechanisms identified in the goal-conditioned setting
(\S\ref{sec:wm-as-env-synthesis}). The diagnosis is unchanged by the
move to single-task: as a source of training-time supervision for the
policy or critic, the LeWM predictor on this benchmark does not work.

\paragraph{What the failures do not rule out.}
None of the four failure modes touches the much weaker claim that the
WM, applied for one or two steps starting from a real on-policy latent,
can produce an informative \emph{ranking} of candidate action chunks
under the learned critic. The WM does not have to be a faithful
long-horizon simulator to do that --- it only has to preserve the
ordering its rollouts induce over a small candidate set well enough
that the highest-scoring action is also a good one in the real
environment. The rest of this chapter exploits exactly that opening.

% ---------------------------------------------------------------------
\section{The pivot: WM at inference only}
\label{sec:owm-pivot}

We continue to train the agent on the offline buffer with real
$(z, \mathbf{a}, r, z')$ transitions, exactly as in
Chapters~\ref{chap:wm-as-env}--\ref{chap:wm-as-subgoal}. The critic
never sees a WM-imagined next state. What changes is how action chunks
are selected at decision time. At each decision step we use the WM as
a rollout-and-score oracle: sample $N$ candidate chunks from the
BC-flow actor at the current latent, integrate each one forward in the
WM for $H$ steps, and score the resulting trajectory under the critic.
The first chunk of the best-scoring candidate is executed. Then, where
this inference-time procedure produces actions that beat the actor's
own samples \emph{under the current critic}, we distil those actions
back into the actor during continued offline training. The supervision
pairs the actor is updated on are $(z, \mathbf{a}^{\text{MPC}})$ where
$z$ is a real latent and only the action label is changed; the critic
is updated, as before, on the real $(z, \mathbf{a}, r, z')$ tuples
from the offline buffer.

% ---------------------------------------------------------------------
\section{Best-of-$N$ MPC at inference}
\label{sec:mpc-bon}

The simplest instantiation samples $N$ candidates from the BC-flow
actor, evaluates each by rolling it forward $H$ steps in the WM and
scoring the resulting trajectory under the critic, and executes the
first chunk of the best candidate. The generic MPC machinery ---
proposal sampling, $H$-step rollouts in $f_\psi$, scoring functions
(Q-term / Q-all / dense$+$Q) --- is described in the background
(\S\ref{sec:mpc-bg}); we focus here on what is specific to the BoN
variant.

\paragraph{Proposal sampling and the connection to Q-chunking.}
Proposals are drawn directly from the BC-flow head of the trained FQL
actor: sample $N$ noise vectors $x_0^{(i)} \sim \mathcal{N}(0, I_{25})$
and integrate the flow with $S = 10$ Euler steps to obtain $N$
candidate action chunks. This is the same best-of-$N$ proposal-selection
scheme that Q-chunking introduces for offline-to-online RL
\cite{li2025reinforcement}; the only addition is to score each
candidate not by the critic at the current state but by the critic
after an $H$-step WM rollout.

\paragraph{Continuation re-sampling.}
When rolling a candidate forward, we draw fresh continuation actions
from the BC-flow actor at the predicted latents rather than reusing the
original candidate's action sequence open-loop. This makes the rollout
an unbiased estimate of the on-policy continuation distribution.

\paragraph{Scoring function.}
The headline numbers in this chapter use \emph{Q-term} scoring --- the
critic evaluated only at the terminal $H$-step state --- because
empirically this is the most reliable: the critic is trained on real
states and is more accurate than the dense reward built on the WM's
geometry, especially in the regimes where the WM's calibration is
sensitive (§\ref{sec:owm-exp2}).

\paragraph{Headline result.}
Algorithm~\ref{alg:mpc-bon} summarises the procedure. On the $81.6\%$
offline checkpoint, the best configuration ($N = 32$, $H = 1$, Q-term)
reaches $88.0\%$ real-env success --- a $+6.4\%$ absolute lift purely
from inference-time search.

\paragraph{Why $H = 1$ wins.}
Increasing $H$ to $2$ or $3$ does not help and often hurts: even one
extra WM step injects enough prediction error that the critic's
evaluation of the terminal state becomes less reliable than its
evaluation at the immediate next state.

\paragraph{The bottleneck.}
Standalone BoN-MPC plateaus around $88$--$90\%$ even with larger $N$,
because the candidate set is bounded by the BC-flow actor's proposal
distribution. To go further we have to either explore actions the
actor would not naturally propose --- which is what
§\ref{sec:mpc-grad}--§\ref{sec:mpc-fmq} do --- or improve the actor
itself --- which is what §\ref{sec:mpc-interleaved} does.

\begin{algorithm}[t]
\caption{Best-of-$N$ MPC inference.}
\label{alg:mpc-bon}
\begin{algorithmic}[1]
\Require current latent $z_t$, BC-flow actor $\pi_\theta$, critic $Q_\theta$, world model $f_\psi$, hyperparameters $N$, $H$, scoring $J$.
\State Tile $z_t$ to a batch $z_{\text{batch}} \in \mathbb{R}^{N \times 192}$.
\State Sample $\{x_0^{(i)}\} \sim \mathcal{N}(0, I_{25})$ and set $a_0^{(i)} \leftarrow \Phi_\theta(z_{\text{batch}}, x_0^{(i)})$.
\State Roll each candidate forward $H$ steps through $f_\psi$ with on-policy continuation re-sampling.
\State Compute score $s_t^{(i)}$ for each candidate using $J$.
\State \Return $a_0^{(i^\star)}$ with $i^\star = \arg\max_i s_t^{(i)}$.
\end{algorithmic}
\end{algorithm}

% ---------------------------------------------------------------------
\section{Gradient-refined MPC}
\label{sec:mpc-grad}

BoN-MPC is restricted to the actions the proposal distribution
samples. Because the WM rollout is differentiable, we can take the
best candidate and refine it further by gradient ascent on the score,
treating the continuation actions as constants so that the gradient
flows only through the initial chunk $a_0$.

\paragraph{Update rule.}
Starting from the BoN winner $a_0$, we run $K$ raw gradient ascent
steps,
\begin{equation}
    a_0 \;\leftarrow\; \mathrm{clip}\!\Big(a_0 + \eta\,\nabla_{a_0} J(a_0;\,z_t),\;-1,\;1\Big),
    \qquad k = 1, \dots, K,
    \label{eq:mpc-grad}
\end{equation}
with step size $\eta = 0.01$ and $K=5$. We refer to this scheme as
\emph{Grad-MPC}; it generalises the standard sampling-MPC controller
of TD-MPC \cite{hansen2022tdmpc, hansen2024tdmpc2} by replacing
resampling with backpropagation through the WM rollout.

\paragraph{Result.}
With $H = 2$, $K = 5$, $N = 32$, Q-term scoring: $90.0\%$ --- a further
$+2.0\%$ over BoN-MPC and now within $2.8\%$ of the IMPALA baseline.

\paragraph{Why Grad-MPC is not the final answer.}
Eq.~\eqref{eq:mpc-grad} is a heuristic: $K$ raw gradient steps with a
fixed step size, with no built-in trust region. With small $K$ the
refinement is too gentle; with large $K$ it overshoots into regions
where the critic has not been trained. The next section replaces this
heuristic with a closed-form update that is, in a precise sense, the
unique optimum.

% ---------------------------------------------------------------------
\section{Flow Map $Q$-Guidance (FMQ)}
\label{sec:mpc-fmq}

The conceptual move is to formulate the action-refinement problem as a
trust-region optimisation problem and ask for the velocity update that
maximises a first-order expansion of the critic under a Euclidean
constraint. The closed-form solution is what \cite{ziakas2026aligning}
call \emph{Flow Map $Q$-Guidance} (FMQ).

\paragraph{Trust-region objective.}
Let $\pi^{\mathrm{ref}}(\cdot \mid z)$ be a reference flow-map policy
producing actions $a_1 = a_r + (1-r)\,u^{\mathrm{ref}}_{r,1}(a_r \mid z)$
through the two-time average velocity $u^{\mathrm{ref}}_{r,1}$
(\S\ref{sec:fmq-bg}). We seek the velocity perturbation that maximises
the first-order expansion of $Q_\phi$ around $a_1$ subject to a
Euclidean ball of radius $\eta$ in velocity space:
\begin{equation}
    \max_{u_{r,1}}\; Q_\phi\!\big(z,\, a_r + (1-r)\, u_{r,1}(a_r \mid z)\big)
    \qquad\text{s.t.}\qquad
    \big\|u_{r,1} - u^{\mathrm{ref}}_{r,1}\big\|_2 \;\le\; \eta.
    \label{eq:mpc-fmq-objective}
\end{equation}
Theorem~3.2 of \cite{ziakas2026aligning} solves this in closed form:
\begin{equation}
    \boxed{\;u^{\,*}_{r,1}(a_r \mid z) \;=\; u^{\mathrm{ref}}_{r,1}(a_r \mid z) \;+\; \eta\,\frac{\nabla_a Q_\phi(z, a_1)}{\big\|\nabla_a Q_\phi(z, a_1)\big\|_2}\;.}
    \label{eq:mpc-fmq-target}
\end{equation}
This is \emph{the} closed-form optimum of
Eq.~\eqref{eq:mpc-fmq-objective} --- any other perturbation within the
$\eta$-ball can only do worse in the local linear approximation.
The magnitude of the update is fixed by $\eta$ rather than by the
magnitude of the gradient; the direction is the normalised gradient.

\paragraph{Inference-time FMQ.}
Used as a planner, FMQ proceeds exactly like Grad-MPC except that the
inner ascent loop of Eq.~\eqref{eq:mpc-grad} is replaced by a single
application of Eq.~\eqref{eq:mpc-fmq-target}. The two methods share the
BoN proposal sampling and the scoring; they differ only in the
refinement step. Algorithm~\ref{alg:mpc-refine} presents both as
branches of the same procedure.

\begin{algorithm}[t]
\caption{Gradient-refined MPC. Grad-MPC ($K$ raw steps) and FMQ (one normalised step) share the proposal-and-score machinery and differ only in the inner refinement.}
\label{alg:mpc-refine}
\begin{algorithmic}[1]
\Require $z_t$, $\pi_\theta$, $Q_\theta$, $f_\psi$, scoring $J$; refinement-specific hyperparameters: $(K, \eta_{\text{Grad}})$ for Grad-MPC, $\eta$ for FMQ.
\State Sample initial proposals $\{a_0^{(i)}\}_{i=1}^N \sim \pi_\theta(\cdot \mid z_t)$.
\State Pre-sample continuation noises held constant during refinement.
\If{Grad-MPC}
    \For{$k = 1, \dots, K$}
        \State $g \leftarrow \nabla_{a_0}\, J(a_0;\,z_t)$ via BPTT through the $H$-step WM rollout.
        \State $a_0 \leftarrow \mathrm{clip}\big(a_0 + \eta_{\text{Grad}}\, g,\;-1,\;1\big)$.
    \EndFor
\ElsIf{FMQ}
    \State $g \leftarrow \nabla_{a_0}\, J(a_0;\,z_t)$ via a single BPTT pass.
    \State $a_0 \leftarrow \mathrm{clip}\big(a_0 + \eta\, g / \|g\|_2,\;-1,\;1\big)$.
\EndIf
\State Draw fresh continuation noises and re-score $a_0^{(1)}, \dots, a_0^{(N)}$ as in Algorithm~\ref{alg:mpc-bon}.
\State \Return $a_0^{(i^\star)}$ with $i^\star = \arg\max_i s_t^{(i)}$.
\end{algorithmic}
\end{algorithm}

\paragraph{Sensitivity on the base critic.}
The FMQ-optimal direction only matters if the critic's value estimate
in that direction is reliable. The critic of the offline FQL agent was
trained on real $(z, a, r, z')$ tuples in which $a$ was sampled from
the BC-flow actor. FMQ-shaped actions --- actions displaced along the
normalised $Q$-gradient --- are not in the distribution the critic has
seen. Empirically, the base critic peaks under FMQ at $\eta \approx 0.1$
($90\%$) and collapses sharply for $\eta > 0.1$ (only $57.2\%$ at
$\eta=0.3$). Large normalised steps move the action into a region of
unreliable $Q$, the score ordering collapses, and FMQ is no better than
random sampling. This is the empirical motivation for
§\ref{sec:mpc-interleaved}: if we train the critic on FMQ-shaped
actions, it learns to be reliable on the action distribution FMQ
produces, and we can use much finer trust regions at inference.

\paragraph{FMQ as a training-time learning target.}
Eq.~\eqref{eq:mpc-fmq-target} also serves a second role beyond
inference. If we set the reference velocity $u^{\mathrm{ref}}_{r,1}$
equal to the frozen offline flow-map velocity $u^{\text{off}}_{r,1}$,
the right-hand side becomes a closed-form self-bootstrapped learning
target for the online flow-map velocity $u^\theta_{r,1}$:
\begin{equation}
    \mathcal{L}_{\text{FMQ}}(\theta) \;=\; \mathbb{E}_{r, a_0, \mathbf{a}_{\text{data}}}\!\left[\,
        \Big\| u^\theta_{r,1}(a_r \mid z) \;-\; \mathrm{sg}\!\Big(u^{\text{off}}_{r,1}(a_r \mid z) \;+\; \eta\,\hat g_\phi\Big) \Big\|_2^2\,\right],
    \quad \hat g_\phi = \frac{\nabla_a Q_\phi(z, a_1)}{\|\nabla_a Q_\phi(z, a_1)\|_2 + \kappa},
    \label{eq:mpc-fmq-loss}
\end{equation}
using a $u$-time interpolant $a_r = (1 - r)\,a_0 + r\,\mathbf{a}_{\text{data}}$
between noise $a_0 \sim \mathcal{N}(0, I)$ and an offline action
$\mathbf{a}_{\text{data}}$. This is Eq.~(12) of the FMQ paper. We use
it as the actor-side label generator inside the interleaved
distillation loop of §\ref{sec:mpc-interleaved}, and it is the source
of the $96\%$ headline result. The training-time use is a contribution
of this thesis; the inference-time use is borrowed from the FMQ paper.

\paragraph{Offline-trained Q as an OOD indicator.}
A useful observation about FMQ is that the inference-time step size
$\eta$ is implicitly an indicator of how far we can trust the WM's
predictions. Where the offline-trained Q-function assigns high values
and low ensemble disagreement, the WM's short-horizon predictions are
reliable enough that fine-grained FMQ refinement ($\eta=0.02$) finds
good actions. Where it assigns low values or high disagreement, the
critic-gradient direction is itself unreliable, and FMQ collapses to
worse-than-baseline. The offline-trained Q acts as an OOD indicator
for the WM's predictive support, without explicitly modelling the
WM's uncertainty --- the critic's calibration on the offline
distribution does the work.

% ---------------------------------------------------------------------
\section{Offline MPC distillation}
\label{sec:mpc-offline-distill}

The simplest way to fold the inference-time gains into the policy is a
two-phase pipeline: relabel the offline dataset with MPC-selected
actions in a one-shot pass, then continue training the BC-flow actor on
those relabels.

\paragraph{Phase A --- relabel.}
For every latent observation $z \in \mathcal{Z}_{\text{off}}$ in the
offline cache we run Grad-MPC ($N = 32$, $H = 2$, $K = 5$,
$\eta = 0.01$, Q-term) using the offline-trained actor and critic and
the fully fine-tuned WM, and record the resulting chunk as
$\mathbf{a}^{\text{MPC}}(z)$.

\paragraph{Phase B --- distil.}
We continue training only the BC-flow actor on the relabelled pairs,
with the critic and one-step distillation heads frozen. The loss is
Eq.~\eqref{eq:flow-bc-clean} with $\mathbf{a}$ replaced by
$\mathbf{a}^{\text{MPC}}$. The optimiser is AdamW with learning rate
$10^{-4}$, weight decay $10^{-4}$, batch size $256$, for $10^5$
gradient steps.

\paragraph{Phase C --- evaluate.}
The distilled actor can be evaluated standalone (BoN $N = 4$ under the
original critic, no WM at test time) or used as the proposal
distribution inside Grad-MPC or FMQ.

\paragraph{The limitation of one-shot distillation.}
Standalone performance does lift modestly --- to $\sim 81$--$84\%$ ---
but no further. The issue is structural: the labels in Phase A are
generated by a fixed actor and a fixed critic, and once they are baked
into the dataset they are stale for the rest of training. As the actor
improves under $\mathcal{L}_{\text{distil}}$, the labels are no longer
optimal under the improved proposal distribution; as the critic could
be improved if it were also being updated, the scoring function used
to generate the labels is also stale. The interleaved variant fixes
both.

% ---------------------------------------------------------------------
\section{Interleaved MPC distillation}
\label{sec:mpc-interleaved}

The MPC labels depend on both the current actor (which proposes
candidates) and the current critic (which scores them). If we let the
actor improve under those labels and the critic improve alongside it
on real offline transitions, the labels themselves become better on
the next refresh --- the actor proposes better candidates and the
critic scores them more reliably. This is a positive-feedback loop;
one-shot distillation discards it; the interleaved variant preserves
it.

\paragraph{Mechanism.}
Inside the standard offline FQL training loop, we maintain a circular
buffer $\mathcal{B}_{\text{MPC}}$ of $(z, \mathbf{a}^{\text{MPC}}(z))$
pairs. After a warm-up of $T_{\text{warm}}$ steps, the buffer is
refreshed every $T_{\text{relabel}}$ gradient steps:
$|\mathcal{B}_{\text{MPC}}|$ states are sampled uniformly from the
offline buffer and the MPC procedure --- either Grad-MPC or FMQ --- is
run on each using the \emph{current} actor and critic. At every
training step we draw a real minibatch and replace the first
$\lfloor pB \rfloor$ rows with samples from $\mathcal{B}_{\text{MPC}}$,
then run the standard FQL update on the resulting mixed batch.

\paragraph{Critically, the critic still sees only real transitions.}
The replacement happens to both the observations and the actions in
the affected rows, but the rewards and next-states stay anchored to
the real offline tuples. The actor is supervised on
$(z, \mathbf{a}^{\text{MPC}})$ where $z$ is a real latent; the critic
is updated with the real reward and next-state stored alongside that
latent in the offline dataset. The WM never appears in the critic's
TD target.

\paragraph{Mixed vs separate injection.}
An alternative scheme runs the MPC-label update as an actor-only step
on top of the main FQL update, leaving the critic's batch composition
untouched. Empirically the mixed mode wins decisively: the critic must
\emph{see} MPC actions to assign them high $Q$; if the critic is
updated only on actor samples while the actor is being pulled toward
MPC labels, the actor is dragged in a direction the critic cannot
endorse and the loop fails to close.

\paragraph{Grad-MPC labels vs FMQ labels.}
The label generator inside Algorithm~\ref{alg:mpc-interleaved} can be
either Grad-MPC or FMQ. Grad-MPC labels produce a strong but not
optimal result: the headline Grad-MPC-distilled run (\emph{Grad-MPC
distil, full, mixed}) reaches $92.4\%$ on MPC-eval. FMQ labels --- the
closed-form optimum of the trust-region problem --- produce
qualitatively cleaner action displacements that the critic adapts to
more cleanly. The FMQ-distilled run \emph{FMQ-distil, slow-relabel}
($T_{\text{relabel}} = 50\,000$) reaches a standalone peak of $86.8\%$
at $400$k steps and, evaluated with FMQ inference at the fine inference
radius $\eta = 0.02$, reaches \textbf{$96.0\%$}, $+3.2\%$ over the
IMPALA baseline.

\paragraph{Optional advantage filter.}
An AWR-style filter that retains only labels for which
$Q_\phi(z, \mathbf{a}^{\text{MPC}}) > Q_\phi(z, \pi_\theta(z))$ is
straightforward to add. Empirically it gives no consistent improvement
over plain mixed injection on this task, so we report it for
completeness but use the plain variant as the default.

\begin{algorithm}[t]
\caption{Interleaved MPC distillation.}
\label{alg:mpc-interleaved}
\begin{algorithmic}[1]
\Require offline dataset $\mathcal{D}_{\text{off}}$, FQL agent $(\pi_\theta, Q_\phi)$, frozen WM $f_\psi$, hyperparameters $T_{\text{warm}}$, $T_{\text{relabel}}$, buffer size $|\mathcal{B}_{\text{MPC}}|$, mix ratio $p$, label generator $\mathcal{M} \in \{\text{Grad-MPC},\, \text{FMQ}\}$.
\State $\mathcal{B}_{\text{MPC}} \leftarrow \emptyset$.
\For{$\text{step} = 1, \dots, T_{\text{total}}$}
    \If{$\text{step} \ge T_{\text{warm}}$ \textbf{and} $\text{step} \bmod T_{\text{relabel}} = 0$}
        \State Sample $|\mathcal{B}_{\text{MPC}}|$ states from $\mathcal{D}_{\text{off}}$.
        \State For each $z$, set $\mathbf{a}^{\text{MPC}}(z) \leftarrow \mathcal{M}(z;\,\pi_\theta,\,Q_\phi,\,f_\psi)$ (Alg.~\ref{alg:mpc-refine}).
        \State $\mathcal{B}_{\text{MPC}} \leftarrow \{(z,\,\mathbf{a}^{\text{MPC}}(z))\}$; reset the actor-only optimiser state.
    \EndIf
    \State Sample minibatch $B$ from $\mathcal{D}_{\text{off}}$.
    \If{$\mathcal{B}_{\text{MPC}} \ne \emptyset$}
        \State Sample $\lfloor pB \rfloor$ pairs from $\mathcal{B}_{\text{MPC}}$ and overwrite the first $\lfloor pB \rfloor$ rows of $B$'s observations and actions.
    \EndIf
    \State Update $(\pi_\theta, Q_\phi)$ on $B$ with the standard FQL critic and actor losses.
\EndFor
\end{algorithmic}
\end{algorithm}

% ---------------------------------------------------------------------
\section{Goal-conditioned validation}
\label{sec:owm-gc-validation}

The bulk of this chapter operates on the single-task variant of the
benchmark because that is where the OGBench IMPALA reference number
sits. We additionally evaluated the final FMQ-distilled checkpoint on
the original goal-conditioned variant (\emph{visual-cube-single-v0},
five tasks $\times$ $50$ episodes), to confirm we have not
inadvertently switched to a method that only works on a single task.
The result is in the same ballpark as the single-task run for the
tasks where the offline policy already had non-trivial competence; the
method does not collapse on goal-conditioned cube but also does not
match the IMPALA reference on every task, since the per-task variance
on the goal-conditioned benchmark is wider than on the single-task
variant. This is the evidence that the method is not switching the
problem definition out from under us --- it is the same agent applied
to a wider goal distribution, not a method that requires single-task
framing to work.

% ---------------------------------------------------------------------
\section{Offline-to-online inside the world model}
\label{sec:mpc-online-wm}

Every method up to this point has been careful to keep the world model
out of the critic's training signal. The two distillation schemes of
\S\S\ref{sec:mpc-offline-distill}--\ref{sec:mpc-interleaved} change only
the \emph{action} the actor is supervised on; the critic continues to
bootstrap on the real offline tuples $(z, \mathbf{a}, r, z')$. This
section asks the question those schemes deliberately avoid: \emph{can we
reintroduce world-model transitions into the replay buffer at all?} This
is exactly the offline-to-online-inside-the-WM setup that
\S\ref{sec:owm-exp1} showed collapses --- the agent acts in the WM, gets
imagined feedback, and trains on it. We show that with one structural
change it not only stops collapsing but recovers the strongest
standalone policy of the chapter, and we explain precisely why.

\paragraph{The hidden inconsistency in action-only distillation.}
The interleaved scheme of \S\ref{sec:mpc-interleaved} overwrites the
action of a real transition with the MPC-selected chunk and leaves the
stored reward and next-state untouched, so the critic is updated on
$(z,\, \mathbf{a}^{\text{MPC}},\, r,\, z')$. We argued there that this
keeps the WM out of the TD target, which is true; but it introduces a
different defect. The next-state $z'$ stored alongside $z$ is the
consequence of the \emph{data} action $\mathbf{a}$, not of
$\mathbf{a}^{\text{MPC}}$. Taking $\mathbf{a}^{\text{MPC}}$ from $z$
would, in general, land somewhere else entirely. The tuple the critic is
asked to fit therefore describes a transition that never happened: it
pairs an action with the reward and successor of a \emph{different}
action. The loop closes --- the critic does learn to assign high $Q$ to
$\mathbf{a}^{\text{MPC}}$, which is why the mixed scheme outperforms the
actor-only one --- but it closes on a dynamically inconsistent target,
and the only thing preventing this from corrupting the value function is
that the displacement $\mathbf{a}^{\text{MPC}} - \mathbf{a}$ is small
enough that $z'$ is still approximately right.

\paragraph{World-model-consistent transitions.}
The principled repair is to stop reusing the stale successor and instead
ask the WM what $\mathbf{a}^{\text{MPC}}$ actually leads to. We run the
agent as a closed-loop controller \emph{inside the world model}: from a
real initial latent we select a chunk by FMQ-MPC (\S\ref{sec:mpc-fmq}),
step the WM to obtain the imagined successor
$z'_\psi = f_\psi(z, \mathbf{a}^{\text{MPC}})$, assign a sparse
goal-reaching reward
$r_\psi = \mathbb{1}\{\lVert z'_\psi - z_{\text{goal}}\rVert < \tau\}$,
and continue the rollout from $z'_\psi$ until success or a $40$-step
truncation. Each step contributes a transition
$(z,\, \mathbf{a}^{\text{MPC}},\, r_\psi,\, z'_\psi)$ that is now
\emph{self-consistent} --- the successor really is the model's
prediction for the action taken --- and these are added to a replay
buffer seeded with the offline data. The cost of consistency is that
$z'_\psi$ is imagined: we have traded the action-mismatch of
\S\ref{sec:mpc-interleaved} for a direct dependence on the quality of
the WM. This is the dependence the WM-as-environment experiments of
\S\S\ref{sec:owm-exp1}, \ref{sec:owm-exp3} and \ref{sec:owm-synthesis}
spent their length warning against.

\paragraph{Two trusts that make the dependence safe.}
The danger is real, but it is bounded by two assumptions, each much
weaker than ``the WM is a faithful simulator,'' and each supported by
the diagnostics earlier in the chapter.
\begin{enumerate}
    \item \textbf{We trust states our value function can rank.} The
    chunk that is rolled out is not an arbitrary action --- it is the
    $\arg\max$ over a candidate set under the \emph{offline} critic,
    which was trained only on real $(z,\mathbf{a},r,z')$ and is
    therefore reliable on the data manifold. The critic acts as an
    in-distribution gate: a candidate whose imagined trajectory wanders
    off the manifold is scored by a $Q$ the critic cannot vouch for and
    is rejected in favour of one that stays on it. The WM is never
    required to be globally accurate, only locally rank-preserving over
    a handful of one- to two-step rollouts --- precisely the weak claim
    left open by \S\ref{sec:owm-synthesis}.
    \item \textbf{We trust the WM to rank a small set of actions.} Using
    the model to choose between a few candidate action chunks over a
    short horizon is exactly how the LeWM architecture uses it in the
    original paper --- its CEM planner (\S\ref{sec:mpc-bg}) at
    evaluation time \cite{maes2026leworldmodelstableendtoendjointembedding}.
    We are not asking the WM to do anything it was not built to do; we
    reuse its sanctioned mode of operation as the data-generation
    mechanism for an online phase.
\end{enumerate}
The combination is what distinguishes this setup from the failure of
\S\ref{sec:owm-exp1}: there, raw WM rollouts entered the buffer with no
value-based gate on which states were admissible, and the critic
bootstrapped on all of them; here, the clean offline critic both
\emph{selects} the actions that generate the data and is kept frozen so
it never has to bootstrap on the imagined successors.

\paragraph{Experimental setup (E2).}
We test this on a single task,
\texttt{visual-cube-single-singletask-task1-v0}, using the fully
fine-tuned LeWM. The starting point is a fresh pure-offline FQL
checkpoint (seed $42$) whose standalone success is $76.4\%$ at the final
step and $88.0\%$ at its peak. The online phase runs $500\,000$ gradient
steps; at each step one FMQ-MPC action (full config: $N = 32$, $H = 2$,
$\eta = 0.02$, Q-all scoring with one FMQ refinement step) is selected
and rolled in the WM, one imagined transition is appended to the buffer,
and a minibatch is drawn from the combined offline-plus-imagined buffer.
Crucially the WM and the planner are used \emph{only} to generate
training data: evaluation is standalone (best-of-$4$ under the critic,
no planner, no WM at test time), $250$ episodes, reported as
mean $\pm$ std over three seeds. The agent never touches the real
environment during the online phase.

\paragraph{The critic must not bootstrap on imagined transitions (E2a vs.\ E2c).}
The two runs in Table~\ref{tab:mpc-online-e2} differ in exactly one bit.
In \textbf{E2a} the standard FQL update is applied to the mixed buffer,
so the critic bootstraps on the imagined transitions; the result is a
collapse to $30.7 \pm 20.8\%$, with individual seeds finishing at $60$,
$14$, and $18\%$ after transient early peaks near $90\%$. Two mechanisms
compound here. The first is the value-corruption mechanism of
\S\ref{sec:owm-exp3} --- the WM in the TD target. The second is a
reward-convention clash: the offline OGBench transitions carry the
native cost-to-goal reward ($r = -1$ per step until success,
$Q \in [-100, 0]$) \cite{ogbench_park2025}, whereas the imagined
transitions carry the goal-reaching reward $r_\psi \in \{0, 1\}$
($Q \in [0, 1]$); a single critic cannot fit two incompatible value
scales in one buffer. In \textbf{E2c} we freeze the critic for the
entire online phase --- only the BC-flow actor is trained, distilled
toward $\mathbf{a}^{\text{MPC}}$ at the WM-visited latents --- so the
clean offline $Q$ remains the ranker and the reward clash is moot (the
only consumer of rewards, the critic, never updates). This single change
lifts the standalone policy to $90.7 \pm 0.8\%$ final and
$94.8 \pm 0.7\%$ peak: a $+14$-point gain over the offline start,
recovered with no real-environment interaction and no planner at test
time. This is the experimental counterpart of the
``critic anchored to real transitions'' claim of
\S\ref{sec:owm-discussion}: E2a violates it and collapses, E2c respects
it and succeeds.

\paragraph{Goal-conditioned extension (E3).}
The same agent extends to a dedicated goal-conditioned controller with
no change to the RL core: the observation becomes the concatenation
$[z, z_{\text{goal}}]$, the agent is trained jointly across all five
\texttt{visual-cube-single} tasks, and we add hindsight relabelling
\cite{her} with the HIQL-style indicator reward
\cite{park2024hiqlofflinegoalconditionedrl} in place of any
latent-distance reward. (This is a separate construction from the
sanity check of \S\ref{sec:owm-gc-validation}, which merely evaluated
the single-task distilled checkpoint on the goal-conditioned benchmark.)
A pure-offline goal-conditioned checkpoint reaches a $5$-task standalone
average of $90.0\%$, and the online-in-WM phase of E2c is then applied
on top. Table~\ref{tab:mpc-online-e3} reports the ablation. The honest
finding is that the online phase does \emph{not} improve over the strong
offline baseline on final success: the best variants peak around $92\%$
but settle at $80$--$85\%$, which we attribute to coverage dilution (the
buffer drifts from fully diverse offline play toward the narrow slice of
latent space the MPC controller visits) compounded by the larger
$(z, z_{\text{goal}})$ space the goal-conditioned policy must cover. The
leakage-free ``fair'' variants, in which online rollout goals are random
achieved states rather than the evaluation goals, are the strongest
(E3g and E3g-clean, both $\approx 84\%$ final). One consistency check is
worth highlighting: because the goal-conditioned setting uses a single
indicator-reward convention everywhere, the reward clash of E2a is
absent, and leaving the critic \emph{unfrozen} (E3e) no longer collapses
--- direct corroboration that the E2a collapse was driven by the
reward-convention mismatch and the WM-in-target mechanism rather than by
the mere presence of imagined data.

\paragraph{Layering the planner back on.}
E2c yields a strong standalone actor; the natural last step is to apply
the test-time FMQ-MPC planner (\S\ref{sec:mpc-fmq}) on top of its best
checkpoints, as the distillation arc did to reach $96.0\%$
(Table~\ref{tab:mpc-fmq-distilled}). Here it does not help: planning on
each seed's best checkpoint yields $93.2 \pm 0.7\%$
(Table~\ref{tab:mpc-online-e2}), \emph{below} the $94.8\%$ standalone
peak of those same checkpoints. Unlike the distillation arc --- where
the standalone actor was a weak $86.8\%$ proposal distribution that
planning substantially improved --- the E2c actor has already absorbed
$500\,000$ steps of MPC supervision, so re-deriving the same search at
test time adds latency without accuracy. The planner failing to help is
itself the cleanest evidence that the online phase has successfully
distilled the controller into the policy.

% ---------------------------------------------------------------------
\section{Discussion}
\label{sec:owm-discussion}

\paragraph{Why this works where WM-as-environment training failed.}
The world model is used at inference only, over short horizons
($H \le 2$), and only as a source of \emph{rankings} between candidate
actions. The critic continues to be trained on real
$(z, \mathbf{a}, r, z')$ from the offline buffer, so the four failure
mechanisms of §\ref{sec:owm-synthesis} --- compounding error,
$Q$-corruption, threshold pathology, model exploitation --- never come
into play. The structural distinction matters: even the
interleaved-distillation step that does feed WM-derived information
into training only changes the actor's BC target on real states; the
critic's calibration is anchored to real transitions throughout.

\paragraph{Training-inference symmetry under FMQ.}
The base critic peaks under FMQ inference at $\eta \approx 0.1$ and
collapses above; the FMQ-distilled critic peaks at $\eta \approx
0.02$--$0.03$ and stays robust across $[0.01, 0.05]$. The shift is
large enough that it is worth naming the mechanism explicitly: the
$Q$-function calibrates to the action distribution it is supervised
on, so training on FMQ-shaped actions makes the critic reliable on
FMQ-shaped actions, which then permits much finer trust-region
refinement at inference.

\paragraph{The world model does not generalise OOD.}
The recurring theme across Chapters~\ref{chap:wm-as-env}--\ref{chap:online-with-wm}
is that the LeWM predictor is reliable on the offline support and
unreliable everywhere else. Goal-conditioned RL pushes the policy off
the support (Chapter~\ref{chap:wm-as-env}). Long-horizon WM rollouts
push the latent off the support (§\ref{sec:owm-exp1}). Joint
WM-and-critic training pushes the WM's own outputs off the support
(§\ref{sec:owm-exp3}). The methods in this chapter succeed because
they never push past the support: short rollouts from a real latent
stay within the WM's reliable horizon, and the offline-trained Q
acts as an OOD indicator that quietly downweights actions the
predictor cannot be trusted to score. The same structural reason
explains why the WGSP attempt of Chapter~\ref{chap:wm-as-subgoal}
failed: its LL distillation loss collapsed the policy onto a single
action, which collapsed the WM rollout diversity, which pushed the
high-level signal off the support. Interleaved MPC distillation
avoids that route by supervising the actor on real-state, MPC-relabelled
\emph{actions}, not on WM-imagined trajectories.

\paragraph{Limitations.}
Most experiments in this chapter are on a single OGBench task; the
exceptions are the goal-conditioned validation of
\S\ref{sec:owm-gc-validation} and the goal-conditioned runs of
\S\ref{sec:mpc-online-wm}, which span all five \texttt{cube-single}
tasks jointly but do not yet improve on the offline baseline. There is
no real online interaction at any point; distillation is offline
throughout, and the ``online'' phase of \S\ref{sec:mpc-online-wm} is
online only with respect to the world model. The dominant compute
cost is the periodic relabelling pass, which scales linearly in $N$,
$H$, and the buffer size; we have not explored amortising this across
tasks. The success of the inference-time procedure depends on the WM
being accurate over short horizons --- a result we verified on the
cube benchmark in \S\ref{sec:wm-drift} but which is task-dependent.
The FMQ paper also describes a complementary inference-time refinement,
$Q$-Guided Beam Search, that combines stochastic renoising with beam
search over the trust region (\S\ref{sec:fmq-bg}); evaluating it on
top of our distilled actors is a natural extension we leave for
future work.
